\subtituloA{El modelo cliente-servidor}

Hoy en día, la arquitectura de los sitios web se basa en el modelo cliente-servidor, que consta de tres elementos comunes esenciales:

\lista{
  \textbf{Servidor:} Equipo donde se instalan los elementos hardware y software, las técnicas incrustadas y los lenguajes de programación. Para implementar estas tecnologías web se aconseja utilizar un lenguaje de programación que soporte multihilos, permitiendo atender múltiples solicitudes simultáneas de manera eficiente. |
  \textbf{Máquina cliente:} Equipo del usuario donde se ejecutan los programas que realizan la petición de solicitud web al servidor a través de los navegadores. Estos navegadores deben soportar las tecnologías del sitio web: HTML, CSS, JavaScript y controles ActiveX. |
  \textbf{Conectividad de red:} Infraestructura y características utilizadas para la transmisión de datos entre cliente y servidor, incluyendo protocolos de comunicación, seguridad y velocidad de conexión.
}

El objetivo fundamental de un navegador es presentar la información, escrita normalmente en lenguaje HTML y alojada en un servidor con el servicio web instalado, interpretando y renderizando el contenido para el usuario final.

\nota{Para el examen}{
¿Dónde tendremos implementados los diferentes programas para realizar peticiones de solicitud en el servidor?

\textbf{Respuesta:} En la máquina cliente, concretamente en el navegador web, que es el programa encargado de realizar las solicitudes al servidor. Además, ¿el funcionamiento de un navegador se encarga de alojar una página web para responder a las distintas solicitudes recibidas?

\textbf{Respuesta:} Falso. El navegador no aloja la página; es el servidor quien la aloja y sirve, mientras que el navegador solicita, descarga y presenta dicha información al usuario.
}
